\chapter{Standards of Specification for Megascale Programs}
\label{ch:16}

The program includes a family of proposals at planetary and astronomical scale, among them continental-scale freshwater transport corridors, systems for supporting ecological recovery, geothermal energy infrastructure, momentum-exchange tethers, deployable and expandable structures in orbit, space-elevator concepts, and stellar-scale energy collection. These proposals span very different levels of technological readiness, and the program's treatment of them depends on keeping the differences in view. A body of theoretical literature on orbital tethers has existed since the 1970s, and tether concepts have been examined in considerable engineering detail, whereas a terrestrial space elevator requires materials with specific strength beyond what has been demonstrated at the necessary scale and faces unresolved problems of dynamics, debris, and construction, and a Dyson-scale collector involves quantities of mass and energy that place its construction at a different order of difficulty again. Presenting all such proposals as equally near-term would mislead, and the program therefore requires that each be classified by readiness, and that any proposal that depends on an unresolved step be marked as such. Three classes are used: projects that are scientifically understood but economically or institutionally difficult, projects that require genuine scientific or engineering breakthroughs, and projects whose feasibility remains conjectural.

For each proposal the program prescribes a standard of specification before evaluation. A proposal for geothermal infrastructure integrated with mass acceleration, for example, is at present a speculative concept, and the requirement is that it supply an explicit mechanism, an energy balance, the geological constraints on the site and the depth, and the material limits of the components. A proposal for continental freshwater corridors must state elevation changes, pumping energy, evaporative and seepage losses, effects on ecological connectivity, and the governance arrangements across jurisdictions, and it must be compared against alternatives that do not require the construction of the corridor. A proposal for forest restoration must distinguish the restoration of an ecosystem from its replacement by an artificial substitute, and it must avoid afforestation that is inappropriate to the local ecology, since planting where a different biome belongs can damage rather than repair. Some proposals use terminology whose technical meaning has not been fixed, in which case the first task is a definition adequate to evaluate the physical function. The common requirement is that a proposal state what would have to be true for it to work and what would show that it does not.

A further consideration concerns the relation of these programs to present obligations. The long time horizon and the very large scale of megascale projects make them vulnerable to a characteristic form of reasoning in which the magnitude of the eventual benefit is invoked to justify expenditure and sacrifice in the present. The program resists this by a sequencing rule: measures intended to satisfy present obligations, such as clean water, energy access, health, and education for existing persons, must depend only on projects of the first readiness class, so that they do not hinge on breakthroughs that may not occur. The rule is a screening condition and not a guarantee, because a project in the first class can still fail for ordinary reasons of financing, governance, construction, opposition, or environmental effect, and the relevant distinction is between dependence on an unresolved scientific or engineering breakthrough and exposure to implementation risk, which every project carries. Programs in the second and third classes may be pursued as research, with resources proportioned to their prospects and subject to independent review, but they may not be made prerequisites for the fulfillment of what is owed now.

The chapters that follow in this part and the next apply the standard to the principal proposals, taking the terrestrial programs first. Each proceeds in the same order: the concept is stated in the author's own terms where a prior statement exists and otherwise by a working definition that is flagged as provisional, the physical relations on which its feasibility depends are written out, the dependencies on other programs and on governance are listed, the readiness class is assigned with reasons, and a calculation is supplied that either supports or undermines the concept's claim to be taken seriously. Two cautions govern the reading of the calculations. A calculation that shows a concept to be consistent with the laws of physics shows only that it has survived the first test, and the concept remains in the second or third readiness class until its other conditions are met. A calculation that shows a concept to be infeasible under its stated assumptions shows only that the assumptions require revision, and does not by itself refute the underlying idea.

\section*{Formalization}

Let $\mathcal{P}$ be a finite set of megascale programs with a readiness function $r:\mathcal{P}\to\{1,2,3\}$ whose values denote, in order, understood but costly, requiring breakthrough, and conjectural. Let $\mathcal{O}$ be a set of present-obligation measures with a dependency relation $\prec\subseteq\mathcal{O}\times\mathcal{P}$, where $o\prec p$ means that measure $o$ depends on program $p$.

\begin{definition}[Readiness-respecting plan]
A plan is \emph{readiness-respecting} if for every $o\in\mathcal{O}$ and every $p$ with $o\prec p$ (including transitive dependence) one has $r(p)=1$.
\end{definition}

\begin{proposition}
In a readiness-respecting plan the fulfillment of every present-obligation measure is independent of the outcome of every program with $r(p)\ge2$. Conversely, if some $o\prec p$ with $r(p)\ge2$, the fulfillment of $o$ inherits the uncertainty of a breakthrough that the success of programs with $r=1$ does not determine.
\end{proposition}

\begin{proof}
In a readiness-respecting plan the programs on which any $o$ depends all have $r=1$, so the fulfillment of $o$ is a function of the outcomes of those programs alone. In the converse case the outcome of $p$ with $r(p)\ge2$ is not determined by the outcomes of programs of class $1$, and $o$ inherits its uncertainty.
\end{proof}

\begin{remark}[Screening, not guarantee]
The proposition is close to definitional, since it follows from how dependence is defined, and it should be read as a design constraint on plans and not as an empirical claim. Let $I_p$ denote the event that program $p$ with $r(p)=1$ is successfully implemented. If fulfillment of $o$ requires every program on which it depends, then $\Pr(o\ \text{fulfilled})\le\Pr\big(\bigcap_{o\prec p}I_p\big)$, which may be well below $1$ because of financing, governance, construction error, political opposition, or unforeseen environmental effects. Readiness-respect is therefore a necessary screening condition for plans whose obligations are not hostage to unresolved breakthroughs, and it is not a sufficient guarantee of delivery. Implementation risk must be analyzed separately for each program.
\end{remark}

Three physical relations illustrate the quantitative character of the specifications; the first is developed in Chapter 18, the second in Chapter 24, and the third in Chapter 25. For the transport of water through a rise of height $h$ with a pipeline of length $L$, diameter $D$, flow speed $v$, and friction factor $f$, the head loss is $h_f=f\,\tfrac{L}{D}\,\tfrac{v^2}{2g}$ (the Darcy--Weisbach relation), and the minimum energy per unit volume of lifting against the total head is
\[
\frac{E}{V}=\frac{\rho g\,(h+h_f)}{\eta},
\]
where $\rho$ is the density of water, $g$ the gravitational acceleration, and $\eta\in(0,1]$ the combined efficiency of pumping. For $h=1000$ m, $h_f=0$, and $\eta=1$ this gives $\rho g h\approx 9.8\times10^6$ J per cubic meter, about $2.7$ kWh per cubic meter, which sets a floor that no design can undercut and against which alternatives such as desalination can be compared on a common basis.

For a tether of constant working stress $\sigma$ and density $\rho_m$ extending from the surface at radius $R$ to geostationary radius $r_g$, define the effective potential $\Psi(r)=-GM/r-\tfrac12\omega^2r^2$, where $G$ is the gravitational constant, $M$ the mass of the Earth, and $\omega$ its rotation rate. Then the cross-sectional area of a tether of uniform stress satisfies $\sigma\,dA/dr=\rho_m A\,\Psi'(r)$, whence
\[
\frac{A(r_g)}{A(R)}=\exp\!\Big[\frac{\rho_m}{\sigma}\big(\Psi(r_g)-\Psi(R)\big)\Big],
\]
and the taper ratio is determined by the specific strength $\sigma/\rho_m$ of the material: because $\Psi(r_g)-\Psi(R)>0$, the ratio grows exponentially as specific strength falls, which is the quantitative reason that readiness for this concept depends on materials that have not been demonstrated. Finally, a collector capturing a fraction $f$ of the solar luminosity $L_\odot\approx3.8\times10^{26}$ W yields power $P=fL_\odot$, so that the scale of such a program is set by the choice of $f$ and by the mass required to intercept it.
